% Fichier complété par tools/complete_missing_corrections.py.
% Source : GEO/GEO-01/18_Maxpid/18_Maxpid.tex
% Statut : rédigé (4 question(s)).
\begin{corrigebox}[Corrigé rédigé]
\CorrigeQuestion{1}
Le graphe comporte un sommet par solide et les
                    liaisons visibles sur le schéma : pivot pour chaque axe de rotation,
                    glissière pour chaque translation et contact ponctuel/sphère-plan pour
                    le poussoir ou le galet. Les arêtes sont annotées par leur centre,
                    leur axe et, pour le roulement, la condition $V_I=0$.
\CorrigeQuestion{2}
On reconstruit la configuration en appliquant les
                rotations dans l'ordre de la chaîne cinématique puis les translations le
                long de leurs axes. Les longueurs des barres restent constantes. Les
                coordonnées finales se calculent avec
                $\vec i'=\cos q\,\vec i+\sin q\,\vec j$; le tracé doit conserver les
                liaisons et le contact sans pénétration.
\CorrigeQuestion{3}
On reconstruit la configuration en appliquant les
                rotations dans l'ordre de la chaîne cinématique puis les translations le
                long de leurs axes. Les longueurs des barres restent constantes. Les
                coordonnées finales se calculent avec
                $\vec i'=\cos q\,\vec i+\sin q\,\vec j$; le tracé doit conserver les
                liaisons et le contact sans pénétration.
\CorrigeQuestion{4}
On reconstruit la configuration en appliquant les
                rotations dans l'ordre de la chaîne cinématique puis les translations le
                long de leurs axes. Les longueurs des barres restent constantes. Les
                coordonnées finales se calculent avec
                $\vec i'=\cos q\,\vec i+\sin q\,\vec j$; le tracé doit conserver les
                liaisons et le contact sans pénétration.
\end{corrigebox}
